\scanpage{31}
Mais la donnée de $(s_0,u)$ équivaut à celle d'un \uline{répère} de $\Pi$, i.e. d'un isom \deleted{entre $\Pi$ et}
\[
\Pi\simeq\Pi_0(n)
\]
(où $n$ \deleted{\unsure{$\in\mathbb N^*$}} est un entier $\geq3$ \uuline{ou $+\infty$}), et la donnée de a) \deleted{\unsure{b)|a)}} équivaut donc à la donnée de (14) (sauf la dernière \deleted{\unsure{relation|condition}} \inserted{\text{condition de (14) si $n=+\infty$}}), \deleted{et} celle de (b) équivaut à la donnée de $s_0\in E$ satisfaisant la \deleted{première} relation $\sigma_{\deleted{\unsure{0}}E}s_0=s_0$, et la condition (19) se traduit par (22). En résumé

\uline{Prop} 5 Les \deleted{repr.} réalisations géom \inserted{\text{\uline{régulières}}} de $\Pi_0(\infty)$ \inserted{\text{dans $E$}} correspondent aux \deleted{couples} triples $(\sigma_E,\sigma'_{E\deleted{\unsure{*}}},s_0)$ de deux autom affines \deleted{et} \inserted{\text{de $E$}}
\[
\deleted{\unsure{\sigma}}\qquad\sigma_E,\sigma'_E\in\mathrm{Aff}(E)
\]
et d'un \deleted{$s_0\in E$, satisfaisant les}
\[
s_0\in E
\]
satisfaisant les conditions
\[
\begin{aligned}
(1^\circ)\quad &\sigma_E^2={\sigma'_E}^2=\mathrm{id}\\
(2^\circ)\quad &\sigma_Es_0=s_0\\
(3^\circ)\quad &\deleted{\unsure{(1,2,3)}}\ \deleted{\unsure{(u_E=\sigma'_E\sigma_E}}\ 
s_0,s_1=u_Es_0,s^{\deleted{\unsure{2}}}_2=u_E^2s_0\quad\text{non alignés}
\end{aligned}
\]
(\deleted{où} \inserted{\text{où on a}} \deleted{\unsure{$u$}}\inserted{\text{posé}}
\[
u_E=\sigma'_E\sigma_E
\]
% archive sheet number: 30
